\documentclass[10pt, a4paper]{article}

% Packages:
\usepackage[
    ignoreheadfoot, % set margins without considering header and footer
    top=9 mm, % seperation between body and page edge from the top
    bottom=9 mm, % seperation between body and page edge from the bottom
    left=11 mm, % seperation between body and page edge from the left
    right=11 mm, % seperation between body and page edge from the right
    footskip=1.0 cm, % seperation between body and footer
    % showframe % for debugging
]{geometry} % for adjusting page geometry
\usepackage{titlesec} % for customizing section titles
\usepackage{tabularx} % for making tables with fixed width columns
\usepackage{array} % tabularx requires this
\usepackage[dvipsnames]{xcolor} % for coloring text
\definecolor{primaryColor}{HTML}{004F90} % define primary color
\definecolor{tagColor}{HTML}{EAF1F7}
\usepackage{enumitem} % for customizing lists
\usepackage{fontawesome5} % for using icons
\usepackage{amsmath} % for math
\usepackage{ragged2e} % Align text
\usepackage[
    pdftitle={Bach-Khoi Vo - CV},
    pdfauthor={Bach-Khoi Vo},
    pdfcreator={LaTeX with RenderCV},
    colorlinks=true,
    urlcolor=primaryColor
]{hyperref} % for links, metadata and bookmarks
\usepackage[pscoord]{eso-pic} % for floating text on the page
\usepackage{calc} % for calculating lengths
\usepackage{bookmark} % for bookmarks
\usepackage{lastpage} % for getting the total number of pages
\usepackage{changepage} % for one column entries (adjustwidth environment)
\usepackage{paracol} % for two and three column entries
\usepackage{ifthen} % for conditional statements
\usepackage{needspace} % for avoiding page brake right after the section title
\usepackage{iftex} % check if engine is pdflatex, xetex or luatex

% Ensure that generate pdf is machine readable/ATS parsable:
\ifPDFTeX
    \input{glyphtounicode}
    \pdfgentounicode=1
    \usepackage[T1]{fontenc}
    \usepackage[utf8]{inputenc}
    \usepackage{lmodern}
\fi

\usepackage{charter}

\renewcommand{\normalsize}{\fontsize{9.2 pt}{11 pt}\selectfont}

% Some settings:
\raggedright
\AtBeginEnvironment{adjustwidth}{\partopsep0pt} % remove space before adjustwidth environment
\pagestyle{empty} % no header or footer
\setcounter{secnumdepth}{0} % no section numbering
\setlength{\parindent}{0pt} % no indentation
\setlength{\topskip}{0pt} % no top skip
\setlength{\columnsep}{0.15cm} % set column seperation
\pagenumbering{gobble} % no page numbering

\titleformat{\section}{\needspace{4\baselineskip}\color{primaryColor}\bfseries\fontsize{11 pt}{13 pt}\selectfont}{}{0pt}{}[\vspace{1pt}\titlerule]

\titlespacing{\section}{
    % left space:
    -1pt
}{
    % top space:
    5 pt
}{
    % bottom space:
    2 pt
} % section title spacing

\renewcommand\labelitemi{$\vcenter{\hbox{\small$\bullet$}}$} % custom bullet points
\renewcommand\labelitemii{$\vcenter{\hbox{\small$\circ$}}$}
\newcommand{\stacktag}[1]{{\setlength{\fboxsep}{1 pt}\colorbox{tagColor}{\fontsize{7.6 pt}{9 pt}\selectfont\color{primaryColor}#1}}}
\newenvironment{highlights}{
    \begin{itemize}[
        topsep=1.5 pt,
        parsep=0pt,
        partopsep=0pt,
        itemsep=1.5 pt,
        leftmargin=12 pt
    ]
}{
    \end{itemize}
} % new environment for highlights


\newenvironment{highlightsforbulletentries}{
    \begin{itemize}[
        topsep=0.10 cm,
        parsep=0.10 cm,
        partopsep=0pt,
        itemsep=0pt,
        leftmargin=10pt
    ]
}{
    \end{itemize}
} % new environment for highlights for bullet entries

\newenvironment{onecolentry}{
    \begin{adjustwidth}{
        0 cm + 0.00001 cm
    }{
        0 cm + 0.00001 cm
    }
}{
    \end{adjustwidth}
} % new environment for one column entries

\newenvironment{twocolentry}[2][]{
    \def\secondColumn{#2}
    \par
}{
    \unskip\hfill\secondColumn\par
} % new environment for two column entries

\newenvironment{threecolentry}[3][]{
    \onecolentry
    \def\thirdColumn{#3}
    \setcolumnwidth{, \fill, 4.5 cm}
    \begin{paracol}{3}
    {\raggedright #2} \switchcolumn
}{
    \switchcolumn \raggedleft \thirdColumn
    \end{paracol}
    \endonecolentry
} % new environment for three column entries

\newenvironment{header}{
    \setlength{\topsep}{0pt}\par\kern\topsep\centering
}{
    \par\kern\topsep
} % new environment for the header

\newcommand{\placelastupdatedtext}{% \placetextbox{<horizontal pos>}{<vertical pos>}{<stuff>}
  \AddToShipoutPictureFG*{% Add <stuff> to current page foreground
    \put(
        \LenToUnit{\paperwidth-2 cm-0 cm+0.05cm},
        \LenToUnit{\paperheight-1.0 cm}
    ){\vtop{{\null}\makebox[0pt][c]{
        \small\color{gray}\textit{Last updated in Jul 2026}\hspace{\widthof{Last updated in Jul 2026}}
    }}}%
  }%
}%

% save the original href command in a new command:
\let\hrefWithoutArrow\href

% new command for external links:

\begin{document}
    \normalsize

    \begin{header}
        {\fontsize{20 pt}{24 pt}\selectfont \textbf{\color{primaryColor}Võ Bách Khôi}}

        \vspace{2 pt}

        {\color{primaryColor}\faMapMarker}\kern 2.5 pt\mbox{Ho Chi Minh City, Vietnam}\quad
        {\color{primaryColor}\faEnvelope}\kern 2.5 pt\mbox{\href{mailto:itskoiwork@gmail.com}{itskoiwork@gmail.com}}\quad
        {\color{primaryColor}\faPhone}\kern 2.5 pt\mbox{+84 944 157 219}\quad
        {\color{primaryColor}\faGlobe}\kern 2.5 pt\mbox{\href{https://itskoi.github.io/}{Portfolio}}\quad
        {\color{primaryColor}\faLinkedin}\kern 2.5 pt\mbox{\href{https://linkedin.com/in/bachkhoivo}{LinkedIn}}
    \end{header}

    % ================ EXPERIENCE ================
    \section{Experience}
        \begin{twocolentry}{Mar 2025 -- Nov 2026}
            \textbf{AI Engineer} @ \href{https://apps.apple.com/vn/app/wao-food-calorie-counter/id6737834587}{Wao}
        \end{twocolentry}
        \textit{AI calorie-tracking app for Vietnamese users. Sole AI engineer on a $\sim$6-person product team.}
        \begin{highlights}
            \item \textbf{Wao Companion (AI health agent, POC)} \stacktag{LangGraph.js} \stacktag{NestJS}\\
            Built an LLM agent that gives personalized answers to nutrition, exercise and app-support questions, based on each user's profile and in-app behavior and grounded in nutritionist-approved guidelines. The agent plans its claims before writing the reply, and output guardrails check that every figure comes from a tool result and every answer cites its guideline.
            \begin{highlights}
                \item Built the eval with the QC, PO and nutrition teams: hand-written test cases, scripted checks that the cited guidelines match the expected ones, and an LLM tone judge validated by human reviewers. The agent correctly chose whether to answer, clarify or decline in \textbf{93\%} of cases. Of graded replies, \textbf{95\%} covered everything asked, \textbf{92\%} were factually correct and \textbf{89\%} were backed by user data or guidelines.
            \end{highlights}
            \item \textbf{Meal suggestions} \stacktag{Python} \stacktag{PyTorch} \stacktag{LightGBM} \stacktag{FastAPI}\\
            Built a system that suggests meals a user is likely to eat and that fit their remaining calories and macros. Two-tower retrieval over meal templates, then learned ranking of what a user will accept and keep, with portion fitting and rule-based safety filters outside the models. Compared DCN-v2 and LightGBM rankers and chose LightGBM, since DCN-v2 showed no significant gain. In offline evaluation on simulated users: \textbf{new-meal recall@10 +59\%} over popularity, \textbf{NDCG@5 +5\%} over frequency ranking and \textbf{accept rate +13\%} over no ranker.
            \item \textbf{Spec-driven development (SDD) workflow} \stacktag{GitHub Spec Kit} \stacktag{OpenSpec} \stacktag{Playwright} \stacktag{Maestro} \stacktag{GitLab}\\
            One of the first two engineers to adopt Wao's spec-driven AI workflow (Sep 2025): Jira ticket $\rightarrow$ AI-written spec $\rightarrow$ task list $\rightarrow$ implementation $\rightarrow$ AI code review. Pushed for explicit task lists after agents lost context and shipped bugs. Proposed agent-driven E2E testing, adopted on web and mobile.
            \item \textbf{Voice Log search ($\sim$20\% of food logs)} \stacktag{React Native} \stacktag{NestJS} \stacktag{Typesense}\\
            Built the voice-logging search pipeline: a cloud LLM API turns a Vietnamese transcript into structured queries (food, brand, alternative names, tags, grams), which search the food and recipe catalog, with results reranked on preparation state (raw, cooked, processed), brand and food group. \textbf{Recall@3 up 30 percentage points} over keyword search alone, measured on hand-labeled queries. Brought the brand-aware ranking to Meal Scan.
        \end{highlights}

        \vspace{3 pt}

        \begin{twocolentry}{Oct 2023 -- Feb 2025}
            \textbf{AI Engineer} @ Uniquify
        \end{twocolentry}
        \begin{highlights}
            \item \textbf{SoC documentation chatbot} \stacktag{LangGraph Python} \stacktag{FastAPI} \stacktag{Vue.js} \stacktag{vLLM}\\
            Led a 4-person team building a RAG chatbot over $\sim$2 GB of chip-design software manuals for the SoC and IP teams, using self-hosted, on-premises LLMs. Cut command lookup from about a minute of searching PDFs to under 10 seconds, with \textbf{95\% answer accuracy} on a hand-built eval set based on the teams' real lookup tasks.
            \item Built AI training modules (CV, NLP, LLMs, MLOps) and ran seminars and workshops to mentor junior engineers.
        \end{highlights}

        \vspace{3 pt}

        \begin{twocolentry}{Apr 2022 -- Sep 2023}
            \textbf{Research Assistant} @ Computational Linguistics Center (CLC), HCMUS
        \end{twocolentry}
        \begin{highlights}
            \item Co-developed a Vietnamese text-to-image model combining diffusion models and PhoBERT (IEEE RIVF'23, first author).
            \item Improved PhoBERT for Vietnamese sentiment analysis with added linguistic features: F1 95.22\% on the revised UIT-VSFC (FAIR'22).
        \end{highlights}

    % ================ EDUCATION ================
    \section{Education}
        \begin{minipage}[t]{0.43\linewidth}
            \textbf{University of Science, VNU-HCM (HCMUS)}\par
            \textit{Bachelor of Science in Computer Science}\par
            \textbf{GPA:} \textbf{3.74}/4.0 (\textbf{8.93}/10)\hfill Aug 2019 -- Nov 2023\par
            \textbf{Awards:}
            \begin{highlights}
                \item Recipient of the \textbf{Encouragement Scholarship} for academic years 2019, 2020, 2021 and 2022
                \item Recognized for \textbf{Outstanding Research Activities} during the academic year 2022--2023
            \end{highlights}
        \end{minipage}\hfill
        \begin{minipage}[t]{0.54\linewidth}
            \textbf{Certifications:}\par
            \begin{tabularx}{\linewidth}{@{}l@{\hspace{4pt}}>{\raggedright\arraybackslash}X@{\hspace{6pt}}r@{}}
                \labelitemi & \textbf{Google Cloud Skills Boost}, \href{https://www.cloudskillsboost.google/public_profiles/06ce9d7f-cc51-478e-989d-d5fbb9b4c1d4}{Member -- Diamond League} & 2022 -- Present \\
                \labelitemi & \textbf{Berkeley RDI}, \href{https://mcusercontent.com/0d89bb5c8066a9533eb98759d/files/215f6572-b834-33b9-d4bf-230306f19a56/llmagentsf24_certificate_no367.pdf}{LLM Agents MOOC -- Mastery Tier Certificate} & Feb 2025 \\
                \labelitemi & \textbf{IIG Vietnam}, TOEIC \textbf{900} & Jun 2023 \\
                \labelitemi & \textbf{Coursera}, \href{https://www.coursera.org/account/accomplishments/specialization/certificate/U53TP5AD3HD3}{IBM AI Engineering Certificate} & Sep 2022 \\
                \labelitemi & \textbf{Coursera}, \href{https://www.credly.com/badges/3d498c27-3706-484b-a81c-ff26b394ea8b?source=linked_in_profile}{IBM Data Science Professional Certificate} & Jun 2021 \\
            \end{tabularx}
        \end{minipage}

    % ================ PUBLICATIONS ================
    \section{Publications}
        \begin{twocolentry}{Dec 2023}
            \textbf{Combining Diffusion Model and PhoBERT for Vietnamese Text-to-Image Generation} | \textit{IEEE RIVF'23}
        \end{twocolentry}
        \textbf{\textit{Bach-Khoi Vo}}, Anh-Dung Ho, An-Vinh Luong, Dinh Dien. \textbf{DOI:} \href{https://doi.org/10.1109/RIVF60135.2023.10471860}{\textit{10.1109/RIVF60135.2023.10471860}}

        \vspace{3 pt}

        \begin{twocolentry}{Dec 2022}
            \textbf{Sentiment Analysis for Vietnamese Language Using PhoBERT Model} | \textit{FAIR'22}
        \end{twocolentry}
        Thanh-Tu Huynh, \textbf{\textit{Bach-Khoi Vo}}, Anh-Dung Ho, Duc-Lung Vu. \textbf{DOI:} \href{https://doi.org/10.15625/vap.2022.0254}{\textit{10.15625/vap.2022.0254}}

    % ================ TECHNOLOGIES ================
    \section{Technologies}
        \begin{highlights}
            \item \textbf{LLM \& Agents:} LangGraph, LangChain, vLLM, Pydantic
            \item \textbf{Search \& ML:} Typesense, PyTorch, LightGBM, scikit-learn, NumPy, Pandas, Polars, CVXPY
            \item \textbf{Backend \& Full-stack:} Python, TypeScript, NestJS, FastAPI, React, Vue.js, GraphQL, SQL
            \item \textbf{Infra:} PostgreSQL, Redis, Docker, Alembic, GCP
            \item \textbf{Dev Environment:} Bash, Git, Vim, tmux, Linux, macOS
        \end{highlights}

\end{document}
